\chapter{Thuật toán song song cho tổng tiền tố}
\label{ch:algorithms}

Chương này trình bày bốn cách tiếp cận song song hóa scan theo thứ tự lịch sử và theo mức độ
``gần phần cứng'': thuật toán Hillis--Steele (\cref{sec:alg-hs}), thuật toán work-efficient của
Blelloch (\cref{sec:alg-blelloch}), scan chia khối dành cho trường hợp \(n \gg p\)
(\cref{sec:alg-blocked}) và chỉ thị \code{scan} của OpenMP (\cref{sec:alg-omp}). Cuối chương
so sánh các thuật toán trên cả phương diện lý thuyết lẫn lưu lượng bộ nhớ
(\cref{sec:alg-compare}). Trong toàn chương, \(\oplus\) là phép cộng và ví dụ minh họa dùng mảng
\(x = [3, 1, 7, 0, 4, 1, 6, 3]\) của Blelloch~\cite{blelloch1990}.

% Kiểu cho các hình mảng
\tikzset{
  cell/.style={draw=husblue,minimum width=0.95cm,minimum height=0.6cm,inner sep=0pt,font=\small},
  hl/.style={cell,fill=husblue!15},
  hl2/.style={cell,fill=orange!25},
  lbl/.style={font=\footnotesize,anchor=east},
  arr/.style={-{Stealth[length=2mm]},thin,gray!70!black},
}
% \drawrow{y}{label}{v0,...,v7}{danh sách chỉ số tô màu (có thể rỗng)}
\newcommand{\drawrow}[4]{%
  \foreach \h in {#4} {\expandafter\gdef\csname hlx\h\endcsname{}}%
  \node[lbl] at (-0.6,#1) {#2};
  \foreach \v [count=\i from 0] in {#3} {
    \ifcsname hlx\i\endcsname
      \node[hl] (r#1c\i) at (\i*1.05,#1) {\v};
    \else
      \node[cell] (r#1c\i) at (\i*1.05,#1) {\v};
    \fi
  }%
  \foreach \h in {#4} {\global\expandafter\let\csname hlx\h\endcsname\undefined}%
}

\section{Thuật toán Hillis--Steele}
\label{sec:alg-hs}

\subsection{Ý tưởng}

Hillis và Steele~\cite{hillis1986} đề xuất thuật toán scan cho máy Connection Machine với hàng
nghìn bộ xử lý. Ý tưởng dựa trên kỹ thuật \emph{nhân đôi khoảng cách} (recursive doubling): ở
bước \(d = 1, 2, \ldots, \lceil\log_2 n\rceil\), mỗi phần tử \(i \ge 2^{d-1}\) cộng thêm giá trị
của phần tử cách nó \(2^{d-1}\) vị trí về bên trái. Sau bước \(d\), phần tử \(i\) chứa tổng của
\(\min(2^d, i+1)\) phần tử kết thúc tại \(i\); sau \(\lceil\log_2 n\rceil\) bước, nó chứa tổng của
tất cả phần tử \(x_0..x_i\), tức là inclusive scan (\cref{fig:hs}).

\begin{figure}[H]
\centering
\begin{tikzpicture}[y=-1.05cm]
  \drawrow{0}{\(x\)}{3,1,7,0,4,1,6,3}{}
  \drawrow{1}{\(d=1\)}{3,4,8,7,4,5,7,9}{1,2,3,4,5,6,7}
  \drawrow{2}{\(d=2\)}{3,4,11,11,12,12,11,14}{2,3,4,5,6,7}
  \drawrow{3}{\(d=3\)}{3,4,11,11,15,16,22,25}{4,5,6,7}
  \foreach \i in {1,...,7} { \pgfmathtruncatemacro{\j}{\i-1}
    \draw[arr] (r0c\j.south) -- (r1c\i.north); }
  \foreach \i in {2,...,7} { \pgfmathtruncatemacro{\j}{\i-2}
    \draw[arr] (r1c\j.south) -- (r2c\i.north); }
  \foreach \i in {4,...,7} { \pgfmathtruncatemacro{\j}{\i-4}
    \draw[arr] (r2c\j.south) -- (r3c\i.north); }
\end{tikzpicture}
\caption{Thuật toán Hillis--Steele trên 8 phần tử: 3 bước, mỗi bước phần tử \(i\) cộng thêm
phần tử \(i - 2^{d-1}\) (mũi tên). Ô tô màu là các phần tử được cập nhật. Hàng cuối là inclusive
scan.}
\label{fig:hs}
\end{figure}

\subsection{Mã giả và vấn đề cài đặt tại chỗ}

Nếu cài đặt tại chỗ (\(x[i] \mathrel{+}= x[i - 2^{d-1}]\)) trên máy không đồng bộ tuyệt đối như
CPU, một luồng có thể đọc \(x[i - 2^{d-1}]\) \emph{sau khi} luồng khác đã cập nhật nó trong cùng
bước -- một \emph{tranh chấp dữ liệu} (data race) cho kết quả sai. Harris và cộng sự gặp đúng vấn
đề này trên GPU khi các warp không chạy đồng thời và giải quyết bằng \emph{bộ đệm kép}
(double buffering)~\cite{harris2007}: đọc từ một mảng, ghi sang mảng kia, rồi hoán đổi vai trò
(\cref{alg:hs}). Giữa hai bước cần một rào chắn để đảm bảo mọi luồng đã ghi xong.

\begin{algorithm}[H]
\caption{Hillis--Steele với bộ đệm kép}
\label{alg:hs}
\KwIn{\(x[0..n-1]\)}
\KwOut{\(y\) = inclusive scan của \(x\)}
\(a \gets x\); cấp phát mảng phụ \(b[0..n-1]\)\;
\For{\(\mathit{off} \gets 1\); \(\mathit{off} < n\); \(\mathit{off} \gets 2\cdot\mathit{off}\)}{
  \ForPar{\(i \gets 0\) \KwTo \(n-1\)}{
    \lIf{\(i \ge \mathit{off}\)}{\(b[i] \gets a[i - \mathit{off}] + a[i]\)}
    \lElse{\(b[i] \gets a[i]\)}
  }
  \emph{barrier}; hoán đổi \(a \leftrightarrow b\)\;
}
\Return \(a\)\;
\end{algorithm}

\subsection{Phân tích}

\begin{itemize}
  \item \textbf{Span:} \(\Span(n) = \lceil\log_2 n\rceil\) bước -- tối ưu.
  \item \textbf{Work:} bước có \(\mathit{off} = 2^{d-1}\) thực hiện \(n - 2^{d-1}\) phép cộng,
        tổng cộng
        \[
          \Work(n) = \sum_{d=1}^{\log_2 n} \bigl(n - 2^{d-1}\bigr) = n\log_2 n - (n - 1)
          = \Oh(n \log n).
        \]
  \item \textbf{Bộ nhớ:} mỗi bước đọc toàn bộ mảng \(a\) và ghi toàn bộ mảng \(b\):
        \(\log_2 n\) lượt quét, mỗi lượt khoảng \(16n\) byte (với int64) -- tức khoảng
        \(16 n \log_2 n\) byte so với \(16n\) byte của bản tuần tự.
\end{itemize}
Thuật toán \textbf{không work-efficient}. Với \(n = 2^{26}\), nó làm nhiều việc hơn bản tuần tự
khoảng 25 lần; trên 12 luồng, kể cả khi song song hóa hoàn hảo, nó vẫn chậm hơn bản tuần tự
khoảng hai lần. Thuật toán này hợp với phần cứng có số bộ xử lý xấp xỉ \(n\) (như một khối
thread GPU xử lý vài trăm phần tử), nhưng không hợp với CPU.

\section{Thuật toán work-efficient của Blelloch}
\label{sec:alg-blelloch}

\subsection{Ý tưởng: hai lượt quét trên cây nhị phân cân bằng}

Blelloch~\cite{blelloch1990} (dựa trên Ladner và Fischer~\cite{ladner1980}) quan sát rằng cây
nhị phân dùng để tính \emph{reduce} (tổng toàn mảng) đã chứa sẵn nhiều tổng riêng phần, và các
tổng này đủ để dựng toàn bộ prefix sum. Thuật toán gồm hai pha, mỗi pha \(\log_2 n\) bước, thực
hiện \emph{tại chỗ} trên mảng có kích thước \(n = 2^k\) (nếu không thì đệm thêm phần tử đơn vị):

\begin{enumerate}
  \item \textbf{Up-sweep (reduce):} đi từ lá lên gốc. Ở tầng \(d\), mỗi nút cha (lưu tại vị trí
        con phải) nhận tổng của hai con. Kết thúc, \(x[n-1]\) chứa tổng toàn mảng, các vị trí
        khác chứa tổng riêng phần của các cây con (\cref{fig:upsweep}).
  \item \textbf{Down-sweep:} gán phần tử đơn vị (0) cho gốc rồi đi từ gốc xuống lá. Ở mỗi nút,
        con trái nhận giá trị của cha; con phải nhận giá trị của cha cộng với giá trị cũ của con
        trái (tổng của cây con trái từ pha up-sweep). Kết thúc, mảng chứa \emph{exclusive} scan
        (\cref{fig:downsweep}).
\end{enumerate}

\begin{algorithm}[H]
\caption{Blelloch scan (exclusive, tại chỗ, \(n = 2^k\))}
\label{alg:blelloch}
\KwIn{\(x[0..n-1]\)}
\KwOut{\(x\) = exclusive scan của đầu vào}
\tcp{Up-sweep}
\For{\(d \gets 0\) \KwTo \(\log_2 n - 1\)}{
  \ForPar{\(k \gets 0\) \KwTo \(n-1\) \textbf{bước} \(2^{d+1}\)}{
    \(x[k + 2^{d+1} - 1] \gets x[k + 2^{d} - 1] + x[k + 2^{d+1} - 1]\)\;
  }
}
\tcp{Down-sweep}
\(x[n-1] \gets 0\)\;
\For{\(d \gets \log_2 n - 1\) \KwSty{downto} \(0\)}{
  \ForPar{\(k \gets 0\) \KwTo \(n-1\) \textbf{bước} \(2^{d+1}\)}{
    \(t \gets x[k + 2^{d} - 1]\)\;
    \(x[k + 2^{d} - 1] \gets x[k + 2^{d+1} - 1]\)\;
    \(x[k + 2^{d+1} - 1] \gets t + x[k + 2^{d+1} - 1]\)\;
  }
}
\end{algorithm}

\begin{figure}[H]
\centering
\begin{tikzpicture}[y=-1.05cm]
  \drawrow{0}{\(x\)}{3,1,7,0,4,1,6,3}{}
  \drawrow{1}{\(d=0\)}{3,4,7,7,4,5,6,9}{1,3,5,7}
  \drawrow{2}{\(d=1\)}{3,4,7,11,4,5,6,14}{3,7}
  \drawrow{3}{\(d=2\)}{3,4,7,11,4,5,6,25}{7}
  \foreach \a/\b in {0/1,2/3,4/5,6/7} {
    \draw[arr] (r0c\a.south) -- (r1c\b.north); \draw[arr] (r0c\b.south) -- (r1c\b.north); }
  \foreach \a/\b in {1/3,5/7} {
    \draw[arr] (r1c\a.south) -- (r2c\b.north); \draw[arr] (r1c\b.south) -- (r2c\b.north); }
  \draw[arr] (r2c3.south) -- (r3c7.north); \draw[arr] (r2c7.south) -- (r3c7.north);
\end{tikzpicture}
\caption{Pha up-sweep: ở tầng \(d\), phần tử \(k + 2^{d+1} - 1\) cộng thêm phần tử
\(k + 2^d - 1\). Sau \(\log_2 n = 3\) bước, \(x[7] = 25\) là tổng toàn mảng.}
\label{fig:upsweep}
\end{figure}

\begin{figure}[H]
\centering
\begin{tikzpicture}[y=-1.05cm]
  \drawrow{0}{gán 0}{3,4,7,11,4,5,6,0}{7}
  \drawrow{1}{\(d=2\)}{3,4,7,0,4,5,6,11}{3,7}
  \drawrow{2}{\(d=1\)}{3,0,7,4,4,11,6,16}{1,3,5,7}
  \drawrow{3}{\(d=0\)}{0,3,4,11,11,15,16,22}{0,1,2,3,4,5,6,7}
  \draw[arr] (r0c7.south) -- (r1c3.north); \draw[arr] (r0c7.south) -- (r1c7.north);
  \draw[arr] (r0c3.south) -- (r1c7.north);
  \foreach \a/\b in {3/1,7/5} {
    \draw[arr] (r1c\a.south) -- (r2c\b.north); \draw[arr] (r1c\a.south) -- (r2c\a.north);
    \draw[arr] (r1c\b.south) -- (r2c\a.north); }
  \foreach \a/\b in {1/0,3/2,5/4,7/6} {
    \draw[arr] (r2c\a.south) -- (r3c\b.north); \draw[arr] (r2c\a.south) -- (r3c\a.north);
    \draw[arr] (r2c\b.south) -- (r3c\a.north); }
\end{tikzpicture}
\caption{Pha down-sweep: con trái nhận giá trị của cha, con phải nhận cha cộng con trái cũ.
Hàng cuối là exclusive scan \([0, 3, 4, 11, 11, 15, 16, 22]\), khớp với ví dụ của
Blelloch~\cite[Hình 1.4]{blelloch1990}.}
\label{fig:downsweep}
\end{figure}

\subsection{Tính đúng đắn}

\begin{theorem}[Blelloch~\cite{blelloch1990}, Định lý 1.1]
Sau pha down-sweep, mỗi nút của cây chứa tổng của tất cả các lá đứng trước nó theo thứ tự duyệt
tiền thứ tự (preorder). Đặc biệt, lá thứ \(i\) chứa \(x_0 + \cdots + x_{i-1}\).
\end{theorem}

\begin{proof}
Quy nạp từ gốc xuống. Gốc không có lá nào đứng trước nên giá trị đúng là phần tử đơn vị 0 --
đúng như được gán. Giả sử nút cha \(v\) có giá trị đúng, tức bằng tổng các lá thuộc vùng \(A\)
(các lá đứng trước \(v\)).
\begin{itemize}
  \item \emph{Con trái} của \(v\) có đúng tập lá đứng trước là \(A\), vì trong duyệt tiền thứ tự
        con trái được thăm ngay sau \(v\). Do đó gán giá trị của cha cho con trái là đúng.
  \item \emph{Con phải} có tập lá đứng trước là \(A \cup B\), với \(B\) là các lá của cây con trái.
        Tổng các lá của \(B\) chính là giá trị con trái sau up-sweep (được giữ trong biến
        \(t\)). Vì vậy con phải nhận giá trị cha \(+\, t\) là đúng.
\end{itemize}
Theo quy nạp, mọi nút, trong đó có các lá, đều có giá trị đúng. Thứ tự tiền thứ tự trên các lá
trùng với thứ tự trái--phải của mảng, nên lá \(i\) chứa \(\sum_{j<i} x_j\).
\end{proof}

\subsection{Phân tích}

\begin{itemize}
  \item \textbf{Work:} up-sweep thực hiện \(n/2 + n/4 + \cdots + 1 = n - 1\) phép cộng; down-sweep
        thực hiện \(n - 1\) phép cộng và \(n - 1\) phép hoán đổi. Tổng \(\Work(n) = \Oh(n)\): thuật
        toán \textbf{work-efficient}~\cite{harris2007}.
  \item \textbf{Span:} \(\Span(n) = 2\log_2 n\) -- gấp đôi Hillis--Steele, vẫn là \(\Oh(\log n)\).
  \item \textbf{Chuyển sang inclusive:} \(y_i = \text{prescan}_i + x_i\), thêm một lượt song song.
\end{itemize}

\paragraph{Nhược điểm trên CPU.} Dù tối ưu về work, thuật toán có ba điểm yếu khi chạy trên
CPU với dữ liệu lớn:
\begin{enumerate}
  \item \textbf{Truy cập bộ nhớ có bước nhảy \(2^{d+1}\):} từ tầng \(d = 3\) trở lên, mỗi phép
        cộng chạm vào một dòng cache khác nhau, chỉ dùng 8/64 byte. Các tầng thấp đọc lại toàn
        bộ mảng nhiều lần, nên tổng lưu lượng bộ nhớ lớn hơn nhiều so với \(16n\) byte.
  \item \textbf{Nhiều rào chắn:} \(2\log_2 n \approx 52\) barrier với \(n = 2^{26}\); ở các tầng
        trên cùng chỉ còn vài phần tử, đa số luồng ngồi chờ.
  \item \textbf{Đệm lên lũy thừa của 2:} với \(n = 2^k + 1\), kích thước làm việc gần gấp đôi.
\end{enumerate}

\section{Scan chia khối cho trường hợp \texorpdfstring{\(n \gg p\)}{n >> p}}
\label{sec:alg-blocked}

\subsection{Ý tưởng}

Cả hai thuật toán trên đều giả định có \(\Oh(n)\) bộ xử lý. Khi chỉ có \(p \ll n\) bộ xử lý,
Blelloch~\cite[Mục 1.2]{blelloch1990} đề xuất chia mảng thành \(p\) khối liên tiếp, mỗi khối
\(\lceil n/p\rceil\) phần tử, và thực hiện ba pha (\cref{fig:blocked}):
\begin{enumerate}
  \item \textbf{Pha 1 -- reduce cục bộ:} mỗi bộ xử lý \(t\) tính \emph{tuần tự} tổng \(s_t\) của
        khối mình.
  \item \textbf{Pha 2 -- scan các tổng khối:} tính exclusive scan của \([s_0, \ldots, s_{p-1}]\),
        được \(o_t = s_0 + \cdots + s_{t-1}\) -- độ lệch (offset) của khối \(t\). Vì \(p\) nhỏ,
        bước này làm tuần tự (\(\Oh(p)\)) hoặc bằng cây (\(\Oh(\log p)\)).
  \item \textbf{Pha 3 -- scan cục bộ:} mỗi bộ xử lý scan \emph{tuần tự} khối của mình, bắt đầu từ
        \(o_t\).
\end{enumerate}
Đây cũng là cách Harris và cộng sự mở rộng scan GPU cho mảng kích thước bất kỳ: scan từng khối,
ghi tổng mỗi khối vào mảng \code{SUMS}, scan \code{SUMS} thành \code{INCR}, rồi cộng đều
\code{INCR[i]} vào khối \(i\)~\cite[Mục 39.2.4]{harris2007}.

\begin{figure}[H]
\centering
\begin{tikzpicture}[y=-1cm,
  bc/.style={draw=husblue,minimum width=0.72cm,minimum height=0.55cm,inner sep=0pt,font=\small}]
  \def\vals{{4,7,1,0,5,2,6,4,8,1,9,5}}
  \def\outs{{4,11,12,12,17,19,25,29,37,38,47,52}}
  \foreach \i in {0,...,11} {
    \pgfmathtruncatemacro{\b}{int(\i/3)}
    \pgfmathsetmacro{\x}{\i*0.78 + \b*0.35}
    \pgfmathtruncatemacro{\v}{\vals[\i]}
    \pgfmathtruncatemacro{\o}{\outs[\i]}
    \ifodd\b \def\fc{orange!20}\else\def\fc{husblue!12}\fi
    \node[bc,fill=\fc] (a\i) at (\x,0) {\v};
    \node[bc,fill=\fc] (c\i) at (\x,3.3) {\o};
  }
  \node[lbl] at (-0.6,0) {\(x\)};
  \node[lbl] at (-0.6,3.3) {\(y\)};
  \foreach \t/\s/\o in {0/12/0,1/7/12,2/18/19,3/15/37} {
    \pgfmathtruncatemacro{\f}{3*\t} \pgfmathtruncatemacro{\l}{3*\t+2}
    \draw[decorate,decoration={brace,mirror,amplitude=4pt}] (a\f.south west) -- (a\l.south east)
      node[midway,below=5pt,font=\footnotesize] {luồng \t: \(s_\t = \s\)};
    \node[font=\footnotesize,husred] at ($(a\f)!0.5!(a\l)+(0,2.15)$) {\(o_\t = \o\)};
  }
  \node[font=\footnotesize,align=left,anchor=west] at (10.2,0) {Pha 1: tổng\\từng khối};
  \node[font=\footnotesize,align=left,anchor=west] at (10.2,2.15) {Pha 2: exclusive scan\\
    \([12,7,18,15] \to [0,12,19,37]\)};
  \node[font=\footnotesize,align=left,anchor=west] at (10.2,3.3) {Pha 3: scan cục bộ\\từ \(o_t\)};
\end{tikzpicture}
\caption{Scan chia khối với \(n = 12\), \(p = 4\) (ví dụ theo~\cite[Hình 1.6]{blelloch1990}, ở
dạng inclusive). Pha 1 và pha 3 hoàn toàn song song; pha 2 chỉ xử lý \(p\) số.}
\label{fig:blocked}
\end{figure}

\subsection{Phân tích}

Theo Blelloch~\cite{blelloch1990}, thời gian trên EREW PRAM là
\begin{equation}
  T_S(n, p) = 2\bigl(\lceil n/p\rceil + \lceil\log_2 p\rceil\bigr) = \Oh(n/p + \log p),
  \label{eq:blocked-time}
\end{equation}
đạt speedup tối ưu khi \(n/p \ge \log p\) -- điều kiện luôn thỏa trên CPU. Work là
\(2n + \Oh(p)\): khoảng gấp đôi bản tuần tự (mỗi phần tử được cộng ở pha 1 và pha 3), vẫn
\(\Oh(n)\). Chỉ có \textbf{hai} rào chắn, và mọi truy cập bộ nhớ đều \textbf{liên tiếp} -- lý
tưởng cho prefetcher.

\paragraph{Lưu lượng bộ nhớ.} Pha 1 đọc \(x\) (\(8n\) byte), pha 3 đọc lại \(x\) và ghi \(y\)
(\(16n\) byte): tổng \(24n\) byte, gấp 1,5 lần mức tối thiểu \(16n\). Với dữ liệu lớn hơn cache,
thời gian bị chặn dưới bởi \(24n/\beta\): ngay cả song song hoàn hảo, scan chia khối cơ bản cũng
chậm hơn kernel copy khoảng 1,5 lần.

\paragraph{Biến thể scan-then-propagate.} Một cách sắp xếp khác: pha 1 scan cục bộ luôn (ghi
\(y\)), pha 3 cộng offset vào \(y\). Cách này đọc \(x\) một lần nhưng đọc--ghi \(y\) hai lần:
\(8n + 8n + 16n = 32n\) byte, tệ hơn. Báo cáo dùng biến thể reduce-then-scan.

\subsection{Cải tiến: chia tile vừa bộ nhớ đệm}
\label{sec:alg-tiled}

Để loại bỏ lượt đọc thừa, nhóm áp dụng ý tưởng \emph{chia tile} (tinh thần giống kỹ thuật mỗi
luồng xử lý nhiều phần tử của Harris và cộng sự~\cite{harris2007}): thay vì \(p\) khối lớn
\(n/p\), mảng được chia thành các \emph{tile}, mỗi tile gồm \(p\) khối nhỏ \(B\) phần tử, sao cho
\(8B\) byte vừa với cache L2 của một nhân (\cref{fig:tiled}). Ba pha được thực hiện lần lượt cho
từng tile; biến \emph{carry} mang tổng các tile trước sang tile sau. Pha 3 đọc lại khối
\emph{ngay sau} pha 1, khi dữ liệu còn trong L2, nên mảng \(x\) chỉ đọc từ DRAM một lần: lưu
lượng giảm từ \(24n\) xuống \(16n\) byte -- đúng mức tối thiểu.

\begin{figure}[H]
\centering
\begin{tikzpicture}[x=1cm,y=1cm]
  \foreach \t in {0,1,2} {
    \foreach \k in {0,1,2,3} {
      \pgfmathsetmacro{\x}{\t*4.3 + \k*1.0}
      \ifodd\k \def\fc{orange!20}\else\def\fc{husblue!12}\fi
      \node[draw=husblue,fill=\fc,minimum width=1cm,minimum height=0.6cm,font=\scriptsize]
        at (\x,0) {T\k};
    }
    \draw[decorate,decoration={brace,amplitude=4pt}] (\t*4.3-0.5,0.4) -- (\t*4.3+3.5,0.4)
      node[midway,above=5pt,font=\footnotesize] {tile \t};
  }
  \node[font=\footnotesize] at (13.3,0) {\(\cdots\)};
  \draw[->,husred,thick] (1.5,-0.5) .. controls (3.5,-1.1) and (4.5,-1.1) .. (5.8,-0.5)
    node[midway,below,font=\footnotesize] {carry};
  \draw[->,husred,thick] (5.8,-0.5) .. controls (7.8,-1.1) and (8.8,-1.1) .. (10.1,-0.5)
    node[midway,below,font=\footnotesize] {carry};
\end{tikzpicture}
\caption{Blocked-opt: mảng chia thành các tile gồm \(p\) khối \(B\) phần tử (ở đây \(p = 4\),
luồng T\(t\) xử lý khối thứ \(t\) của mỗi tile). Ba pha thực hiện trong từng tile; tổng của tile
được cộng dồn vào \emph{carry} cho tile kế tiếp.}
\label{fig:tiled}
\end{figure}

Cái giá phải trả là số rào chắn tăng từ 2 lên \(\lceil n/(pB)\rceil\) (một rào chắn mỗi tile).
\(B\) vì vậy là tham số đánh đổi: \(B\) nhỏ thì nhiều rào chắn, \(B\) lớn thì khối tràn khỏi L2.
Ảnh hưởng của \(B\) được khảo sát ở \cref{sec:exp-block}.

\section{Chỉ thị \texttt{scan} của OpenMP}
\label{sec:alg-omp}

Với chỉ thị \code{reduction(inscan,+)} (\cref{lst:omp-scan-intro}), lập trình viên không chọn
thuật toán. GCC/libgomp cài đặt chỉ thị này theo sơ đồ chia khối hai lượt: mỗi luồng chạy pha
``cập nhật'' trên đoạn của mình, lưu giá trị cục bộ của biến thu gọn cho \emph{từng} lần lặp vào
bộ đệm tạm, sau đó các tổng đoạn được scan và mỗi luồng chạy pha ``sử dụng'' với giá trị đã hiệu
chỉnh. Bộ đệm tạm có kích thước \(\Oh(n)\) làm tăng lưu lượng bộ nhớ, và thân vòng lặp bị tách
thành hai bản sao nên khó vector hóa. Ưu điểm là ngắn gọn, di động; nhược điểm là hiệu năng phụ
thuộc hoàn toàn vào cài đặt của trình biên dịch.

\section{So sánh các thuật toán}
\label{sec:alg-compare}

\begin{table}[H]
\centering
\caption{So sánh lý thuyết các thuật toán scan (\(n\) phần tử int64, \(p\) luồng, \(B\) kích
thước khối). Lưu lượng DRAM tính cho dữ liệu lớn hơn cache.}
\label{tab:alg-compare}
\renewcommand{\arraystretch}{1.25}
\small
\begin{tabular}{lccccc}
\toprule
\textbf{Thuật toán} & \textbf{Work} & \textbf{Span} & \textbf{Barrier} &
\textbf{Bộ nhớ phụ} & \textbf{Lưu lượng DRAM} \\
\midrule
Tuần tự         & \(n\)          & \(n\)            & 0                & 0          & \(16n\) \\
Hillis--Steele  & \(n\log_2 n\)  & \(\log_2 n\)     & \(\log_2 n\)     & \(n\)      & \(\approx 16n\log_2 n\) \\
Blelloch        & \(2n\)         & \(2\log_2 n\)    & \(2\log_2 n\)    & \(\le 2n\) & \(\gg 16n\) (bước nhảy) \\
Blocked         & \(2n\)         & \(2n/p + p\)     & 2                & \(p\)      & \(24n\) \\
Blocked-opt     & \(2n\)         & \(2n/p + n/B\)   & \(n/(pB)\)       & \(2p\)     & \(16n\) \\
\code{omp scan} & \(\approx 2n\) & \(\approx 2n/p\) & tùy cài đặt      & \(\Oh(n)\) & \(> 24n\) \\
\bottomrule
\end{tabular}
\end{table}

\Cref{tab:alg-compare} gợi ý trước kết quả thực nghiệm: trên CPU, work và span không đủ để dự đoán
hiệu năng; \emph{lưu lượng bộ nhớ} mới là đại lượng quyết định khi dữ liệu lớn. Hillis--Steele và
Blelloch -- hai thuật toán có span \(\Oh(\log n)\) đẹp nhất -- lại có lưu lượng lớn nhất. Ngược
lại, Blocked-opt có span ``kém'' nhất nhưng lưu lượng tối thiểu.
